Expressways are critical infrastructures in modern transportation systems, where traffic flow is dense, vehicle speed is high, and safety margins are relatively small. In such scenarios, low-visibility weather conditions such as fog and localized fog patches can seriously degrade drivers’ perception, delay reaction time, and increase the probability of rear-end collisions and secondary accidents. Compared with urban fog, highway fog clusters are often characterized by strong locality, rapid onset, non-uniform spatial distribution, and severe safety consequences. Therefore, conventional manual monitoring and threshold-based warning strategies are insufficient for real-time and fine-grained perception in intelligent transportation systems. To address this challenge, this thesis proposes a highway fog monitoring method based on monocular depth estimation and unified multi-task learning, aiming to jointly accomplish vehicle detection, fog-type classification, and fog-density regression.

The proposed method starts from practical constraints in engineering deployment. Since large-scale real foggy traffic datasets with dense labels are difficult to acquire, this work reuses clear-weather traffic videos and combines them with precomputed monocular depth and online physically inspired fog synthesis. Specifically, UA-DETRAC is adopted as the base traffic monitoring dataset, and the train/validation split is performed at the video-sequence level to avoid temporal leakage. \code{MiDaS_small} is used to estimate scene-relative depth offline and save depth maps as cache files. During training, a GPU-based online fog augmentation module synthesizes three weather conditions, namely clear, uniform fog, and patchy fog, according to the atmospheric scattering model. This process provides the model not only with detection supervision but also with fog-category labels and a regression target for the scattering coefficient beta.

On top of this data pipeline, a unified multi-task model is designed. The model uses a YOLO detector as the shared backbone and neck, while two lightweight heads are attached to high-level semantic features: a fog classification head and a beta regression head. In this way, target detection, weather recognition, and fog-intensity estimation are jointly optimized in a single network. The system further supports checkpoint recovery, missing-depth-cache auto-completion, mixed-precision training, non-finite gradient monitoring, video inference, adaptive confidence thresholding, and export paths for ONNX and quantization-aware deployment.

Beyond the basic inference chain, the current system also introduces a trajectory-level temporal filtering module to suppress persistent static false positives that visually resemble vehicles in roadside surveillance videos. The filter performs cross-frame association, static-candidate analysis, and optional second-stage appearance verification, and is further connected with a hard-negative mining workflow that exports suspicious tracks, builds manual-review packs, and trains a lightweight binary vehicle/non-vehicle patch classifier. The main contribution of this work is therefore threefold. First, it provides a practical data construction strategy for fog monitoring without relying on large-scale real fog annotations. Second, it explores a unified perception framework for jointly modeling traffic targets and fog attributes. Third, it delivers an end-to-end engineering prototype from training to deployment-oriented inference, including video-level false-positive mitigation and iterative refinement hooks. The resulting system is capable of simultaneously estimating vehicle instances and fog conditions from traffic videos, and can serve as a reproducible baseline for subsequent studies on highway low-visibility warning, traffic meteorology, and edge-side intelligent sensing.
